\subsection{Nonconvex convergence, PL rates, and coupled Lyapunov bounds}
We first analyze a predictable positive preconditioner with an explicit
discrepancy term. This separates the descent mechanism from the costs of
momentum, normalization, clipping, and transport. For STRATA the decomposition
is exact, but obtaining a vanishing error requires control of those costs.
The later analysis of the actual normalization field provides a second route
that retains its dependence on the current gradient directly.

\begin{assumption}[Smooth objective and predictable surrogate]\label{ass:rigid-descent}
For $F:\RR^D\to\RR$, $w_t\in\mathcal F_{t-1}$, and a self-adjoint $B_t$, assume
\begin{gather}
 F_{\inf}:=\inf_wF(w)>-\infty,\quad\EE|F(w_0)|<\infty,\\\quad
 \|\nabla F(w)-\nabla F(v)\|\le L\|w-v\|,\\
 h_t:=\nabla F(w_t),\quad g_t=h_t+b_t+\zeta_t,\quad
 \EE[\zeta_t\mid\mathcal F_{t-1}]=0,\quad
 \EE[\|\zeta_t\|^2\mid\mathcal F_{t-1}]\le\sigma_t^2,\\
 B_t\in\mathcal F_{t-1},\quad 0<mI\preceq B_t\preceq MI,\\
 w_{t+1}=w_t-\eta_t(B_tg_t+e_t),\quad
 \varepsilon_t:=e_t+B_tb_t.
\end{gather}
\end{assumption}
\begin{theorem}[Variance-resolved finite-time stationarity]\label{thm:tf-convergence}
Under Assumption~\ref{ass:rigid-descent}, let
$0<\eta_t\le1/(4LM)$, $S_T:=\sum_{t<T}\eta_t$, and
$\Delta_0:=\EE F(w_0)-F_{\inf}$. Define
\begin{equation}
 \Sigma_t:=\EE_t[\zeta_t\zeta_t^\top],\quad
 \mathsf v_t:=\tr(B_t^2\Sigma_t)=\EE_t\|B_t\zeta_t\|^2
                         \le M^2\sigma_t^2.
\end{equation}
Write $\|x\|_A^2:=x^\top Ax$ for $A\succ0$. Then
\begin{align}
 \EE F(w_{t+1})
 &\le\EE F(w_t)-\tfrac12\eta_t\EE\|h_t\|_{B_t}^2
       +L\eta_t^2\EE\mathsf v_t
       +\eta_t\EE\|\varepsilon_t\|_{B_t^{-1}}^2
       +L\eta_t^2\EE\|\varepsilon_t\|^2,\\
 \frac{\sum_{t<T}\eta_t\EE\|h_t\|^2}{S_T}
 &\le\frac{2\Delta_0}{mS_T}
 +\frac{2L\sum_{t<T}\eta_t^2\EE\mathsf v_t}{mS_T}
 +\frac{2\sum_{t<T}\left[\eta_t\EE\|\varepsilon_t\|_{B_t^{-1}}^2
                         +L\eta_t^2\EE\|\varepsilon_t\|^2\right]}{mS_T}.
 \label{eq:tf-stationarity}
\end{align}
For an independent output index $J$ with $\PP(J=t)=\eta_t/S_T$, the left side is
$\EE\|\nabla F(w_J)\|^2$.
\end{theorem}
\begin{proof}
Use the descent lemma~\citep{nesterov2004introductory} and the positive-operator
inequality $B_t^2\preceq MB_t$. Weighted Fenchel--Young bounds the defect
in the matching dual metric. With $\overline g_t:=h_t+\zeta_t$ and
$Q_t:=\langle h_t,B_th_t\rangle$,
\begin{align}
 F(w_{t+1})&\le F(w_t)-\eta_t\langle h_t,B_t\overline g_t+\varepsilon_t\rangle
                    +\tfrac L2\eta_t^2\|B_t\overline g_t+\varepsilon_t\|^2,\\
 \EE_t\langle h_t,B_t\overline g_t\rangle
 &=\langle h_t,B_th_t\rangle,\\
 |\langle h_t,\varepsilon_t\rangle|
 &\le\tfrac14 Q_t+\|\varepsilon_t\|_{B_t^{-1}}^2,\\
 \EE_t\|B_t\overline g_t+\varepsilon_t\|^2
 &\le2\|B_th_t\|^2+2\mathsf v_t+2\EE_t\|\varepsilon_t\|^2,\\
 \|B_th_t\|^2&\le M\langle h_t,B_th_t\rangle,\\
 -\eta_t(1-LM\eta_t)Q_t
                 +\tfrac14\eta_t Q_t
 &\le-\tfrac12\eta_t Q_t,\\
 \tfrac12\sum_{t<T}\eta_t\EE\|h_t\|_{B_t}^2
 &\le\EE F(w_0)-\EE F(w_T)
 +L\sum_{t<T}\eta_t^2\EE\mathsf v_t\\
 &\quad+\sum_{t<T}\left[\eta_t\EE\|\varepsilon_t\|_{B_t^{-1}}^2
                              +L\eta_t^2\EE\|\varepsilon_t\|^2\right].
\end{align}
Finally $\|h_t\|_{B_t}^2\ge m\|h_t\|^2$. The defect energy is kept directional;
replacing it by $m^{-1}\|\varepsilon_t\|^2$ is an optional later relaxation.
\end{proof}
\begin{corollary}[Sharp descent in the zero-discrepancy case]\label{cor:sharp-descent}
If $\varepsilon_t=0$ and $0<\eta_t<2/(LM)$, then
\begin{equation}
 \sum_{t<T}\eta_t(1-LM\eta_t/2)\EE\|h_t\|_{B_t}^2
 \le\Delta_0+\frac L2\sum_{t<T}\eta_t^2\EE\mathsf v_t.
\end{equation}
\end{corollary}
\begin{proof}
Conditional centering gives the exact second-moment identity
$\EE_t\|B_t(h_t+\zeta_t)\|^2=\|B_th_t\|^2+\mathsf v_t$.
Apply the descent lemma without a two-term norm relaxation and sum.
For $F(w)=L\|w\|^2/2$ and $B_t=MI$, every one-step inequality is an equality:
both the noise coefficient $L/2$ and the endpoint $2/(LM)$ are sharp for this
class. The positive-operator step also explains the larger step range in
Theorem~\ref{thm:tf-convergence}.
\end{proof}
The stationarity bound has three sources: the initial objective gap, accumulated
stochastic variance, and accumulated update discrepancy. Its left side concerns
a step-size-weighted output iterate. Thus a decreasing training loss alone
does not verify the conclusion, and persistent discrepancy can leave a
nonvanishing stationarity floor even as the stochastic term decreases.

\begin{corollary}[Rates and almost-sure stationarity along a subsequence]
If $\EE\mathsf v_t\le\overline v$, $\eta_t=\eta/\sqrt T\le1/(4LM)$, then
\begin{align}
 \EE\|\nabla F(w_J)\|^2
 &\le\frac{2\Delta_0}{m\eta\sqrt T}
       +\frac{2L\eta\overline v}{m\sqrt T}
       +\frac{2}{m}\left(\frac1m+\frac{L\eta}{\sqrt T}\right)
                        \frac1T\sum_{t<T}\EE\|\varepsilon_t\|^2.
\end{align}
If almost surely
\begin{gather}
 \sum_t\eta_t=\infty,\quad
 \sum_t\eta_t^2\mathsf v_t<\infty,\quad
 \sum_t(\eta_t/m+L\eta_t^2)\EE_t\|\varepsilon_t\|^2<\infty,
\end{gather}
then
\begin{equation}
 F(w_t)\text{ converges a.s.},\qquad
 \sum_t\eta_t\|h_t\|^2<\infty\text{ a.s.},\qquad
 \liminf_{t\to\infty}\|h_t\|=0\text{ a.s.}
\end{equation}
\end{corollary}
\begin{proof}
The finite-time assertion follows by substitution; one may always take
$\overline v=M^2\sigma^2$ when $\sigma_t^2\le\sigma^2$, and
$\|\varepsilon_t\|_{B_t^{-1}}^2\le m^{-1}\|\varepsilon_t\|^2$.
For the almost-sure assertion,
apply the nonnegative almost-supermartingale theorem \citep{robbins1971convergence}
to $F(w_t)-F_{\inf}$, the summable conditional error, and
$\tfrac m2\eta_t\|h_t\|^2$; an eventually positive lower bound on $\|h_t\|$
would contradict $\sum_t\eta_t=\infty$.
\end{proof}
\begin{theorem}[PL contraction and an explicit floor]\label{thm:tf-PL}
In addition, suppose
\begin{gather}
 \tfrac12\|\nabla F(w)\|^2\ge\mu(F(w)-F^\star),\quad\mu>0,\\
 0<\eta\le\min\{1/(4LM),1/(\mu m)\},\quad
 \EE\mathsf v_t\le\overline v,\quad\EE\|\varepsilon_t\|^2\le E^2.
\end{gather}
For $a_t:=\EE[F(w_t)-F^\star]$, $q:=1-\mu m\eta$,
\begin{align}
 a_{t+1}&\le qa_t+L\eta^2\overline v+(\eta/m+L\eta^2)E^2,\\
 a_T&\le q^Ta_0+
 \frac{L\eta\overline v+(m^{-1}+L\eta)E^2}{\mu m}(1-q^T).
\end{align}
\end{theorem}
\begin{proof}
\begin{gather}
 -\tfrac m2\eta\EE\|h_t\|^2\le-\mu m\eta a_t,\qquad
 \sum_{j=0}^{T-1}q^j=(1-q^T)/(\mu m\eta).
\end{gather}
\end{proof}
Under PL, the same descent estimate becomes a contraction of objective error.
The formula distinguishes a transient geometric term from the error floor:
reducing the learning rate suppresses the variance contribution, while a fixed
discrepancy budget $E^2$ can persist. The diminishing-step result specifies when
decaying discrepancy also permits the objective error to vanish.

\begin{lemma}[Diminishing-step PL recurrence]\label{lem:rigid-polynomial}
If $t_0>a>\max\{1,p\}$, $p>0$, and
\begin{equation}
 a_{t+1}\le\left(1-\frac a{t+t_0}\right)a_t+
                  \frac b{(t+t_0)^2}+\frac c{(t+t_0)^{1+p}},
\end{equation}
then
\begin{align}
 a_T&\le\left(\frac{t_0}{T+t_0}\right)^a a_0
 +\frac{b\,2^a}{(T+t_0)^a}\sum_{t<T}(t+t_0)^{a-2}
 +\frac{c\,2^a}{(T+t_0)^a}\sum_{t<T}(t+t_0)^{a-p-1}\\
 &=O((T+t_0)^{-a}+(T+t_0)^{-1}+(T+t_0)^{-p}).
\end{align}
Thus $\eta_t=\eta_0/(t+t_0)$, $\EE\|\varepsilon_t\|^2=O((t+t_0)^{-p})$,
and $a=\mu m\eta_0>\max\{1,p\}$ imply the displayed rate.
\end{lemma}
\begin{proof}
\begin{align}
 \prod_{j=s+1}^{T-1}\left(1-\frac a{j+t_0}\right)
 &\le\exp\left(-a\sum_{j=s+1}^{T-1}\frac1{j+t_0}\right)
 \le\left(\frac{s+1+t_0}{T+t_0}\right)^a,\\
 (s+1+t_0)^a&\le2^a(s+t_0)^a,\qquad
 \sum_{t<T}(t+t_0)^u=O((T+t_0)^{u+1})\quad(u>-1).
\end{align}
\end{proof}
\begin{theorem}[Coupled parameter--flag Lyapunov inequality]\label{thm:rigid-lyapunov}
Suppose nonnegative sequences $E_t,R_t,\Xi_t,\sigma_t^2$ satisfy
\begin{align}
 \EE F(w_{t+1})&\le\EE F(w_t)-a\eta\EE\|h_t\|^2
               +c\eta E_t+v\eta^2\sigma_t^2+r\eta R_t,\\
 E_{t+1}&\le\kappa E_t+b\eta^2\EE\|h_t\|^2+q\eta^2\sigma_t^2+s\Xi_t,
 \quad0\le\kappa<1.
\end{align}
Let $a,c>0$, $b\ge0$,
\begin{gather}
 \eta^2\le\frac{a(1-\kappa)}{4bc}\quad(bc=0:\ \text{no constraint}),\quad
 \Lambda:=\frac{2c\eta}{1-\kappa},\quad
 \mathcal L_t:=\EE F(w_t)-F_{\inf}+\Lambda E_t.
\end{gather}
Then
\begin{align}
 \mathcal L_{t+1}-\mathcal L_t
 &\le-\tfrac a2\eta\EE\|h_t\|^2-c\eta E_t
       +(v+\Lambda q)\eta^2\sigma_t^2+r\eta R_t+\Lambda s\Xi_t,\\
 \frac1T\sum_{t<T}\EE\|h_t\|^2
 &\le\frac{2\mathcal L_0}{a\eta T}
 +\frac{2(v+\Lambda q)\eta}{aT}\sum_{t<T}\sigma_t^2
 +\frac{2r}{aT}\sum_{t<T}R_t
 +\frac{4cs}{a(1-\kappa)T}\sum_{t<T}\Xi_t.
\end{align}
\end{theorem}
\begin{proof}
\begin{align}
 \mathcal L_{t+1}-\mathcal L_t
 &\le-(a\eta-\Lambda b\eta^2)\EE\|h_t\|^2
       +(c\eta-\Lambda(1-\kappa))E_t
       +(v+\Lambda q)\eta^2\sigma_t^2+r\eta R_t+\Lambda s\Xi_t,\\
 a\eta-\Lambda b\eta^2&\ge a\eta/2,\qquad
 c\eta-\Lambda(1-\kappa)=-c\eta.
\end{align}
Summation and $\mathcal L_T\ge0$ give the conclusion.
\end{proof}
The joint potential accounts for feedback between parameter motion and flag
tracking: each can introduce error into the other. Its weight $\Lambda$ is
chosen so that contraction of the flag error absorbs its contribution to
parameter descent. The following corollary then connects covariance estimation
and tracking to the stationarity bound through explicit intermediate constants.

\begin{corollary}[Covariance-to-stationarity error chain]\label{cor:tf-endtoend}
Suppose a linear block-isotropic surrogate satisfies
\begin{gather}
 \EE\|\varepsilon_t\|^2\le
   2G_\star^2(\Gamma^a)^2J_P^2 E_t+2R_t,\quad
 d_P(\psi(x),\psi(y))\le J_P\|x-y\|,\\
 E_{t+1}\le\kappa E_t+B_t.
\end{gather}
For constant $\eta\le1/(4LM)$ and $\sup_t\EE\mathsf v_t\le\overline v$,
\begin{align}
 \frac1T\sum_{t<T}\EE\|h_t\|^2
 &\le\frac{2\Delta_0}{m\eta T}+\frac{2L\eta\overline v}{m}\\
 &\quad+\frac4m(m^{-1}+L\eta)\left[
 \frac{G_\star^2(\Gamma^a)^2J_P^2}{T(1-\kappa)}
             \left(E_0+\sum_{t<T-1}B_t\right)+\frac1T\sum_{t<T}R_t\right].
\end{align}
For the implemented coordinatewise optimizer, $R_t$ additionally contains the
normalization, gauge, momentum, clipping, and decay terms of
Definitions~\ref{def:rigid-map} and~\ref{def:rigid-error} and
Lemma~\ref{lem:rigid-clipping}.
\end{corollary}
\begin{proof}
\begin{gather}
 \|(B(U_t)-B(U_t^\star))g_t\|^2
 \le G_\star^2(\Gamma^a)^2J_P^2\|x_t-x_t^\star\|^2,\quad
 \sum_{t<T}E_t\le(E_0+\sum_{t<T-1}B_t)/(1-\kappa).
\end{gather}
\end{proof}
\subsection{A virtual-iterate theorem for matrix momentum}
The virtual iterate subtracts the parameter displacement stored in momentum.
Its increment resembles a preconditioned gradient step, with a remainder that
records changing coefficients and transport. This algebraic reformulation
allows the proof to retain the whole momentum history through variation
budgets, rather than discarding it as an unspecified per-step error.

\begin{assumption}[Predictable momentum system]\label{ass:rigid-momentum}
Let $\mathcal B_t,B_t\in\mathcal F_{t-1}$, with
$\|\mathcal B_t\|_{\op}\le\beta<1$, $mI\preceq B_t\preceq MI$, and
\begin{gather}
 m_t=\mathcal B_t m_{t-1}+(I-\mathcal B_t)g_t+u_t,\quad
 w_{t+1}=w_t-\eta_tB_tm_t,\quad m_{-1}=0,\\
 \EE_tg_t=\nabla F(w_t),\quad\EE_t\|g_t\|^2\le G_2^2,\quad
 \|m_t\|\le G_m,\quad\|\nabla F(z_t)\|\le G_z,\\
 C_t:=\eta_tB_t(I-\mathcal B_t)^{-1}\mathcal B_t,\quad C_{-1}:=0,\quad
 z_t:=w_t-C_{t-1}m_{t-1}.
\end{gather}
\end{assumption}
\begin{lemma}[Exact virtual-iterate identity]\label{lem:rigid-virtual}
Under Assumption~\ref{ass:rigid-momentum},
\begin{gather}
 z_{t+1}=z_t-\eta_tB_tg_t+r_t,\\
 r_t:=(C_{t-1}-C_t)m_{t-1}-\eta_tB_t(I-\mathcal B_t)^{-1}u_t,\\
 \|r_t\|\le G_m\|C_t-C_{t-1}\|_{\op}
                      +\frac{\eta_tM}{1-\beta}\|u_t\|,\qquad
 \|w_t-z_t\|\le\|C_{t-1}\|_{\op}G_m.
\end{gather}
\end{lemma}
\begin{proof}
\begin{align}
 \eta_tB_t+C_t&=\eta_tB_t(I-\mathcal B_t)^{-1},\\
 (\eta_tB_t+C_t)m_t
 &=C_tm_{t-1}+\eta_tB_tg_t
                       +\eta_tB_t(I-\mathcal B_t)^{-1}u_t,\\
 z_{t+1}-z_t&=-(\eta_tB_t+C_t)m_t+C_{t-1}m_{t-1}.
\end{align}
\end{proof}
\begin{theorem}[Stationarity with explicit variation and transport budgets]
\label{thm:rigid-momentum}
Under Assumptions~\ref{ass:rigid-descent} (smoothness and lower bound only)
and~\ref{ass:rigid-momentum},
\begin{align}
 \frac m2\sum_{t<T}\eta_t\EE\|\nabla F(w_t)\|^2
 &\le F(w_0)-F_{\inf}+LM^2G_2^2\sum_{t<T}\eta_t^2\\
 &\quad+\frac{M^2L^2G_m^2}{2m}\sum_{t<T}\eta_t\|C_{t-1}\|_{\op}^2
       +G_z\sum_{t<T}\EE\|r_t\|+L\sum_{t<T}\EE\|r_t\|^2.
\end{align}
For constant scalar momentum $\mathcal B_t=\beta I$ and constant
$\eta=\eta_0/\sqrt T$, define
$\mathcal V_T:=\sum_{t<T}\|B_t-B_{t-1}\|_{\op}$, with $B_{-1}=0$. Then
\begin{align}
 \sum_{t<T}\EE\|r_t\|
 &\le\frac{\eta\beta G_m}{1-\beta}\EE\mathcal V_T
                  +\frac{\eta M}{1-\beta}\sum_{t<T}\EE\|u_t\|,\\
 \EE\mathcal V_T+\sum_{t<T}\EE\|u_t\|=O(\sqrt T),\quad
 \sum_{t<T}\EE\|r_t\|^2=O(1)
 &\Longrightarrow\frac1T\sum_{t<T}\EE\|\nabla F(w_t)\|^2=O(T^{-1/2}).
\end{align}
\end{theorem}
\begin{proof}
Let $h_t=\nabla F(w_t)$, $\delta_t=\nabla F(z_t)-h_t$. Then
\begin{align}
 \|\delta_t\|&\le L\|C_{t-1}\|_{\op}G_m,\\
 -\EE_t\langle\nabla F(z_t),B_tg_t\rangle
 &=-\langle h_t+\delta_t,B_th_t\rangle\\
 &\le-\tfrac m2\|h_t\|^2+\tfrac{M^2}{2m}\|\delta_t\|^2,\\
 F(z_{t+1})&\le F(z_t)-\eta_t\langle\nabla F(z_t),B_tg_t\rangle
 +\langle\nabla F(z_t),r_t\rangle
 +\tfrac L2\|-\eta_tB_tg_t+r_t\|^2,\\
 \EE_tF(z_{t+1})&\le F(z_t)-\tfrac m2\eta_t\|h_t\|^2
 +\frac{\eta_tM^2L^2G_m^2}{2m}\|C_{t-1}\|_{\op}^2
 \notag\\
 &\quad+G_z\EE_t\|r_t\|+LM^2\eta_t^2G_2^2+L\EE_t\|r_t\|^2.
\end{align}
Summation gives the first assertion. Also
$\|C_t\|_{\op}\le\eta M\beta/(1-\beta)$ and
$\sum_t\eta\|C_{t-1}\|_{\op}^2=O(T\eta^3)$.
\end{proof}
